\begin{statement}{A.22}{Proved}{Fixed subgroups on covers with pro-prime Galois closure}
Let $K=\mathbf Q(i)$, $U=E\setminus\{O\}$, and use the chosen
geometric base point, cusp inertia $I$, and packet section $s_b$
of Part A. Write
\[
 \pi=\pi_1(U_{\overline K}),\qquad
 P_\ell=\pi^{(\ell)},\qquad
 I_\ell=\operatorname{im}(I\to P_\ell).
\]
Let $H\subset\pi$ be open, and suppose the finite quotient by
its normal core is an $\ell$-group. Then the natural continuous
map $H^{(\ell)}\to P_\ell$ is injective with open image
$\overline H=\operatorname{im}(H\to P_\ell)$.
If $\Delta_0\subset\Delta$ is any open subgroup such that
$s_b(\Delta_0)$ preserves $H$, then, in this embedding,
\[
       (H^{(\ell)})^{s_b(\Delta_0)}
           =\overline H\cap I_\ell
           =\operatorname{im}(H\cap I\to H^{(\ell)}).
                                                        \tag{A.22.1}
\]
The last group is the inertia of the cusp selected by the coset
$H\cdot1$ and the chosen base-point path.

In particular this computes the full fixed subgroup at $\ell=2$
for all three $D_8$ genus-three covers, the $Q_8$ genus-three cover,
and every genus-two reflection quotient of a $D_8$ cover in A.4.
If $c$ generates the original $I_2$, their fixed subgroup maps
to $\overline{\langle c^2\rangle}$ in $P_2$; intrinsically it is
the full pro-$2$ inertia of the selected cusp. These statements
hold for every packet section, including those with vanishing
degree-two Kummer character, and for every stated $\Delta_0$.
\end{statement}
\begin{proof}
We first verify the finite-index restriction of the fixed-group
input R1 (LM8 G.10). The following elementary observation applies
to any finite-rank free pro-$\ell$ group $P$. If an automorphism
$\alpha$ acts trivially on $P^{\mathrm{ab}}$ and
$\alpha^N(u)=u$ for a positive integer $N$, then $\alpha(u)=u$.
Indeed $\alpha$ acts trivially on every lower-central graded
piece, because these are generated from degree one by brackets.
They are torsion-free $\mathbf Z_\ell$ modules, by the integral
Magnus realization of the free Lie algebra. If
$v=\alpha(u)u^{-1}\ne1$, choose its first nonzero lower-central
class $v_d$. Modulo the next term, the product giving
$\alpha^N(u)u^{-1}$ has class $Nv_d$, which cannot vanish.
The lower-central filtration is separated: every finite
$\ell$-group quotient is nilpotent. This proves the observation,
also when $\ell$ divides $N$.

CM over $K$ makes the action of $\Delta$ on
$P_\ell^{\mathrm{ab}}=T_\ell E$ trivial. If $u$ is fixed by an
open subgroup $\Delta_0$, every $g\in\Delta$ has a positive
power in $\Delta_0$ (one may first take its finite-index normal
core). Apply the observation to the action of $s_b(g)$.
It follows that
\[
           P_\ell^{s_b(\Delta_0)}
             =P_\ell^{s_b(\Delta)}=I_\ell.
                                                        \tag{A.22.2}
\]
The final equality is the all-prime packet statement R1,
whose attached proof LM8 G.10 separates the genuine degree-two
inner case from the nonzero higher metabelian-detected
derivation case. The finite-index step just proved needs
neither a nonzero degree-two coefficient nor a normal
$\Delta_0$; it does not assume that fixed points commute with
pro-$\ell$ completion.

We next prove the claimed completion comparison. Put
$C=\bigcap_{g\in\pi}gHg^{-1}$, so $\pi/C$ is a finite
$\ell$-group. Given an open normal subgroup $N\triangleleft H$
with $H/N$ a finite $\ell$-group, $N\cap C$ is normal in $C$.
Intersect its finitely many conjugates under $\pi/C$ to obtain
$M\triangleleft\pi$, with $M\subset N$.
The quotient $C/M$ embeds in a finite product of finite
$\ell$-groups $C/g(N\cap C)g^{-1}$. Hence both $C/M$ and
$\pi/M$ are finite $\ell$-groups. This proves that the
pro-$\ell$ topology of $H$ is induced from that of $\pi$.
Equivalently, writing $R_\ell(\pi)$ for the kernel of
$\pi\to P_\ell$, one has
\[
            R_\ell(H)=R_\ell(\pi)\subset C\subset H.
\]
Thus $H^{(\ell)}=H/R_\ell(\pi)$ identifies with
$\overline H\subset P_\ell$, and its index is $[\pi:H]$.
The identification is equivariant for every automorphism
preserving $H$, in particular for $s_b(\Delta_0)$.
Intersecting (A.22.2) with $\overline H$ now gives the first
equality of (A.22.1).

Since $H$ is the inverse image of $\overline H$ under
$\pi\to P_\ell$, every element of $\overline H\cap I_\ell$
has a lift in $H\cap I$. The converse inclusion is immediate.
The geometric covering interpretation of the latter subgroup
is precisely the inertia of the selected cusp. This proves
the second equality and identifies the actual subgroup, not
only its conjugacy class after a further quotient.

For the listed applications the Galois closure is $D_8$ or
$Q_8$, so it is a $2$-group. Their puncture monodromy, the
commutator of a generating pair, is the nontrivial central
element of order two. It is absent from a reflection subgroup
of $D_8$, so it has orbits of length two also on the four
cosets defining each genus-two quotient. Thus the ramification
index of every displayed cusp over $O$ is exactly two.
The original inertia generator therefore enters each $H$
first as $c^2$, proving the final formula. The independent
finite certificate \texttt{code/a\_fixed\_subcovers.py}
checks the commutators of all generating pairs in these two
groups and the reflection-coset orbit lengths.
\end{proof}
